\chapter{Directed probability enclosures}\label{ch:intervals}

\label{sec:final-4}

An analytic remainder estimate controls the error of a mathematical
approximation. A numerical enclosure must also control the arithmetic used
to evaluate it. This chapter joins directed arithmetic to explicit
remainder bounds, yielding an interval that contains the probability
represented by an exactly supplied positive spectrum, subject to the stated
arithmetic trust base.

The notation follows the analytic development: $\kappa$ is the cumulant generating function, $b$ the support endpoint, $N$ the retained mode count, and $J_k$ the payoff coefficient. The implementation-name correspondence is given in Appendix~\ref{app:reproduction}.

\section{Interval target}

For a positive central quadratic form
\[
Q=\sum_{j=1}^{r}\lambda_jZ_j^2,
\qquad
\lambda_j>0,
\]
the output is not merely a point estimate.  It is
\[
\boxed{
\log_{10}P(Q\ge q)\in[d_-,d_+].
}
\]
Hence the corresponding multiplicative uncertainty satisfies
\[
\boxed{
\frac{p_+}{p_-}\le 10^{d_+-d_-}.
}
\]

The finite transform and exponential prefactor are evaluated by
directed-rounding interval arithmetic.  The resulting interval is enlarged by
two independently proved analytic remainders: a support/periodization alias
bound and an omitted-COS-mode bound.

\section{No exact saddlepoint is required}

Choose any
\[
0<\theta<\frac{1}{2\lambda_{\max}}.
\]
Exponential tilting gives
\[
P(Q\ge q)
=
e^{\kappa(\theta)-\theta q}
S_\theta(q),
\]
where
\[
S_\theta(q)
=
E_\theta\!\left[
e^{-\theta(Q-q)}
1_{\{Q\ge q\}}
\right].
\]

The interval proof does not require \(\theta\) to be the exact saddlepoint.
A high-precision saddlepoint is used only as an efficient design choice.
Once serialized as an exact decimal input, it is simply a valid tilt inside
the MGF domain.

The same is true of the support endpoint \(b\) and the auxiliary Chernoff
parameter \(\eta\).  Correctness requires only
\[
b>q\ge0,\qquad b>0,
\qquad
\theta<\eta<\frac{1}{2\lambda_{\max}}.
\]

\section{Rigorous support and alias control}

Under the tilted measure,
\[
P_\theta(Q>b)
\le
\exp\left\{
\kappa(\eta)-\kappa(\theta)-(\eta-\theta)b
\right\}.
\]

The full-characteristic-function cosine expansion periodizes the tilted density
on \([0,b]\).  Because the original positive quadratic form is supported on
\([0,\infty)\), the total positive-tail mass aliased back into one fundamental
interval is at most \(P_\theta(Q>b)\).

The actual payoff omitted above \(b\) is at most
\[
e^{-\theta(b-q)}P_\theta(Q>b).
\]

Therefore
\[
\boxed{
E_{\rm support}
\le
P_\theta(Q>b)
\left(
1+e^{-\theta(b-q)}
\right).
}
\]

\section{Rigorous infinite-mode control}

Define
\[
a_j=
\frac{2\lambda_j}{1-2\theta\lambda_j}.
\]
The normalized tilted characteristic function satisfies
\[
|\phi_\theta(t)|
=
\prod_{j=1}^r
\left(
1+(a_jt)^2
\right)^{-1/4}.
\]

Since
\[
1+(a_jt)^2\ge(a_jt)^2,
\]
we have
\[
|\phi_\theta(t)|
\le
C_\theta t^{-r/2},
\qquad
C_\theta=\prod_j a_j^{-1/2}.
\]

For the payoff coefficient
\[
J_k=
\int_q^b
e^{-\theta(x-q)}
\cos(k\pi x/b)\,dx,
\]
the elementary bound
\[
|J_k|\le\frac1\theta
\]
holds.

Let
\[
\nu=\frac r2>1,\qquad N\ge1.
\]
For a finite transform retaining modes \(0,\ldots,N-1\),
\[
\boxed{
E_{\rm series}
\le
\frac{2C_\theta}{b\theta}
\left(\frac{b}{\pi}\right)^\nu
\left[
N^{-\nu}
+
\frac{N^{1-\nu}}{\nu-1}
\right].
}
\]

This first interval backend therefore requires \(r>2\).  Lower-rank positive
forms can use a different tail bound and are not claimed by this theorem.

\section{Final enclosure theorem}

Let \(\boldsymbol S_N\) be the directed-rounding interval containing the finite
\(N\)-mode tilted COS sum, and set
\[
E=E_{\rm support}+E_{\rm series}.
\]
Then
\[
\boxed{
S_\theta(q)
\in
\boldsymbol S_N+[-E,E].
}
\]

If the lower endpoint of this interval is positive, directed rounding gives
\[
\boxed{
\log P(Q\ge q)
\in
\kappa(\theta)-\theta q
+
\log\left(
\boldsymbol S_N+[-E,E]
\right).
}
\]

All finite arithmetic in this expression is enclosed by the interval backend.

\section{Completing the running calculation}\label{sec:worked-enclosure}
The factor, spectrum, and tilt of Sections~\ref{sec:worked-law}--\ref{sec:worked-tilt} now determine a complete positive-spectrum calculation at $q=40$. The exact diagonalization supplies the coefficients $(1,1,2,2)$. The implementation receives those coefficients as exact decimal literals, together with a serialized decimal approximation to the admissible tilt. Correctness does not require that decimal to equal the radical saddlepoint exactly.

With $b=512$, $N=16384$, auxiliary tilt $\eta=0.245$, and 60 interval decimal digits, the supplied backend returns a logarithmic interval. Rounding its lower endpoint downward and its upper endpoint upward gives the readable enclosure
\[
\log_{10}P(Q\ge40)\in\bigl[-4.04489575,-4.03897381\bigr].
\]
At $N=8192$, the directed ratio upper bound is approximately $0.0246131$, which fails the requested tolerance $0.02$. Doubling the retained mode count, rather than relaxing the tolerance, gives the interval above. The full serialized interval has a directed multiplicative ratio upper bound below $0.013731$. The displayed, coarser interval also meets the requested tolerance $0.02$. An independently evaluated interval for the exact expression $2e^{-10}-e^{-20}$ lies entirely inside the COS enclosure.

This same fixed enclosure does not meet tolerance $0.001$. The acceptance routine rejects that tighter request. This is an intended distinction between an accurate point estimate and an established enclosure width. The rank-four example has slower Fourier decay than the rank-eighty control developed later, so its mode count should not be interpreted as a general cost prediction.


The certificate in this example concerns the tail evaluation for the exact spectrum. Its connection to the matrix is the explicit algebraic diagonalization above. It does not test a black-box eigensolver. Chapter~\ref{ch:spectrum} gives the separate post-verification theorem required for a general computed eigensystem.

\section{Independent exact interval control through \texorpdfstring{\(10^{-5000}\)}{1e-5000}}

For the control
\[
Q=0.02\chi^2_{80},
\]
the survival probability has the finite exact formula
\[
P(Q\ge q)
=
e^{-x}
\sum_{j=0}^{39}
\frac{x^j}{j!},
\qquad
x=\frac{q}{0.04}.
\]

This finite expression was evaluated independently by directed-rounding
interval arithmetic, without using K-COS.  The K-COS enclosure and this exact
control enclosure are compared at thresholds targeting
\[
p\approx10^{-100},
\qquad
10^{-1000},
\qquad
10^{-5000}.
\]

With $N=128$ and 90 interval decimal digits, the resulting K-COS widths
are listed below. Widths and ratios are computed from the full outward
endpoints before rounding the displayed values. At all three stored inputs,
the independent exact-law interval is fully contained in the K-COS interval.

\begin{center}
\begin{tabular}{@{}rrr@{}}
\toprule
target & width in $\log_{10}p$ & multiplicative ratio minus one\\
\midrule
$10^{-100}$ &
$8.93\times10^{-13}$ &
$2.06\times10^{-12}$\\
$10^{-1000}$ &
$7.01\times10^{-12}$ &
$1.62\times10^{-11}$\\
$10^{-5000}$ &
$3.34\times10^{-11}$ &
$7.70\times10^{-11}$\\
\bottomrule
\end{tabular}
\end{center}

Thus the \(10^{-5000}\) result is not merely a point estimate agreeing with a
reference.  For the exact supplied spectrum, it comes with a directed-rounding
logarithmic enclosure whose probability ratio satisfies
\[
\boxed{
\frac{p_+}{p_-}
\le
1+7.70\times10^{-11}
}
\]
The displayed upper bound is rounded upward from the directed ratio bound.

\begin{figure}[htbp]
\centering
\includegraphics[width=.82\linewidth]{figure_24_interval_deep_tail_width.png}
\caption{Width of the directed-rounding K-COS probability enclosure on exact
scaled-chi-square controls, using $N=128$ and 90 interval decimal digits.
Widths are recomputed from the full outward endpoints of the retained inputs.
The enclosure remains narrow even near $10^{-5000}$.}
\label{fig:interval-tail-width}
\end{figure}

\section{Arithmetic width versus analytic width}

At the \(10^{-5000}\) control, the interval precision was swept from 40 to 120
decimal digits.  The total certified width remained essentially
\[
3.34247\times10^{-11}
\]
decades.

Over the same sweep, the arithmetic-only width contracted from approximately
\[
7.5\times10^{-37}
\]
to
\[
6.3\times10^{-117}
\]
decades.

\begin{figure}[htbp]
\centering
\includegraphics[width=.82\linewidth]{figure_25_interval_arithmetic_vs_analytic.png}
\caption{Arithmetic versus analytic uncertainty at a $10^{-5000}$ tail.  The
directed-rounding arithmetic enclosure becomes negligible. The final width is
controlled by the explicit analytic support and series remainders.}
\label{fig:interval-arithmetic}
\end{figure}

This separation is important.  Increasing numerical precision cannot improve
an analytic truncation theorem.  Conversely, once interval arithmetic is far
below the analytic error budget, numerical rounding is no longer the limiting
uncertainty.

\section{Guarantee boundary}

The directed enclosure has the following scope:

\begin{quote}
For an exact supplied positive decimal spectrum of rank greater than two, the
backend returns a directed-rounding logarithmic probability enclosure enlarged
by proved analytic support/alias and omitted-mode bounds.
\end{quote}

The result is conditional on correctness of the interval-arithmetic library
implementation.  The implementation uses \texttt{mpmath.iv}, whose
interval layer applies downward and upward endpoint rounding.

A completely end-to-end guarantee beginning from an arbitrary floating-point
matrix requires one more layer: either a rigorously enclosed matrix-to-spectrum
computation, or an exact/residual completion step whose linear algebra is also
performed with verified interval or ball arithmetic.

The packetized spectral representation has likewise not yet been
ported to interval arithmetic.  The first rigorous backend deliberately uses
the exact supplied spectrum so that representation error cannot be hidden
inside the numerical enclosure.



